\section*{الفصل \ref{ch:b3:facies-models} --- نماذج السحنات}

\begin{solution}{exo:b3:facies-models:1}
من البحر نحو اليابسة: المنطقة البحرية المفتوحة، والانتقال البحري المفتوح مع رمال العواصف،
\omterm{def:b3:facies-models:shoreface}{الوجه الشاطئي}، والشاطئ، وسهل ساحلي تنتشر فيه المستنقعات. وقعت كل بيئة فوق جارتها الواقعة نحو البحر،
فتحرّك خط الشاطئ نحو البحر، وهذا انحسار. لو حدث تجاوز بحري لتراكمت
الطبقات بالترتيب المعاكس، فيكون السهل الساحلي عند القاعدة والطين البحري عند القمة،
وغالبًا مع سطح تعرية عند الموضع الذي
مرّ فيه \omterm{def:b3:facies-models:shoreface}{الوجه الشاطئي}.
\end{solution}

\begin{solution}{exo:b3:facies-models:2}
(أ) ردم قناة أو حاجز نقطي: تآكل حصوي عند القاعدة، وتطبّقات متقاطعة،
ثم تموّجات صغيرة وطين به جذور نباتية. (ب) وجه شاطئي متقدّم أو جبهة دلتا:
طين تتخلّله طبقات رمل عاصفي وتطبّق متقاطع رابي في الأسفل،
ورمل \omterm{def:b3:facies-models:shoreface}{الوجه الشاطئي} ذو التطبّق المتقاطع، وشاطئ في الأعلى. (ج) رواسب متراكبة لنهر مضفور
رملية، أو كثيب ريحي، أو حاجز كربوناتي: مجموعات تطبّق متقاطع ذات حجم حبيبي ثابت
مع واجهات كثبانية كبيرة للحالة الريحية، أو حبيبات هيكلية وأوليتية
في حالة الحاجز.
\end{solution}

\begin{solution}{exo:b3:facies-models:3}
$\phi = (2.65 - 2.40)/(2.65 - 1.0) = 0.152$؛ وبالماء الملحي،
$0.25/1.55 = 0.161$.
\end{solution}

\begin{solution}{exo:b3:facies-models:4}
$e^{-2\pi z/\lambda} = e^{-\pi} = 0.043$ عند $z = \lambda/2$، أي نحو
\qty{38}{m} و\qty{55}{m}. عند \qty{20}{m}: $e^{-2\pi\times20/75} = 0.19$
و$e^{-2\pi\times20/110} = 0.32$ من حركة السطح.
\end{solution}

\begin{solution}{exo:b3:facies-models:5}
$(60 - 47.6)/(189 - 47.6) = 0.088$؛ $(60 - 52.6)/(189 - 52.6) = 0.054$. وقيمة
المصفوفة غير المؤكَّدة بمقدار \qty{5}{\micro\second/ft} تغيّر المسامية بمقدار الثلث:
معادلة المتوسط الزمني تحتاج إلى قيمة للمصفوفة تُتحقَّق منها على اللبّ.
\end{solution}

\begin{solution}{exo:b3:facies-models:6}
$(6.0 + 3.5 + 8.0 + 2.5)/40 = 20/40 = 0.50$. حجم المسام: $20\times0.22
\times10^6 = \qty{4.4e6}{m^3}$ لكل كيلومتر مربّع.
\end{solution}

\begin{solution}{exo:b3:facies-models:7}
قرب الكانيون: Ta وحده، إذ أزال الجريانُ التالي الأقسامَ العليا؛ وفي
أسفل المروحة: Tc (تموّجات) و Td (طمي) تحت Te (طين). ويفقد التيار سرعته
وحمولته الخشنة مع المسافة؛ ولا ينتقل إلا الجزء الناعم
\qty{30}{km}، ويترسّب بفعل جريان أبطأ يُنشئ التموّجات، لا
بالتساقط من معلّق كثيف.
\end{solution}

\begin{solution}{exo:b3:facies-models:8}
حاجز نقطي لنهر متعرّج: رواسب متخلّفة على قاع القناة المتآكل،
ورمل ذو تطبّق متقاطع من القناة العميقة، ورمل متموّج في قمة الحاجز،
وطين السهل الفيضي به جذور. ويسجّل التعاقب القناةَ من القاع
إلى الضفة، لذا كان عمق القناة نحو \qty{9}{m} (أقلّ قليلًا، بعد
انضغاط الطين). ولا تمتدّ القراءة إلى ما تحت القاعدة: فهي سطح تعرية،
وما يقع تحتها قد يكون أيّ رسوب أقدم.
\end{solution}

\begin{solution}{exo:b3:facies-models:9}
$60/40\,000 = 0.0015$، \ang{0.09}: منحدر. $300/1500 = 0.20$، \ang{11.3}،
بعد قمة مستوية عمقها \qty{5}{m}: \omterm{def:b3:facies-models:ramp}{رصيف ذو حافة}، مع حبيبات شعابية أو حاجزية
عند الهامش، وكتل وجريانات حبيبية على المنحدر أدناه.
\end{solution}

\begin{solution}{exo:b3:facies-models:10}
يستغرق الحزام زمنًا مقداره $w/v$ لعبور موضع ثابت، وخلاله يرتفع القاع
بالنسبة إلى مستوى سطح البحر بمعدّل $A$ أسفله: يملأ الرمل $h = A\,w/v$.
$2000\times0.2/20 = \qty{20}{m}$. وعند $A = 0$ لا يترك \omterm{def:b3:facies-models:shoreface}{الوجه الشاطئي} إلا
سمكَ مقطعه الخاص، وهو طبقة رقيقة، إذ تنزلق الأحزمة فوق بعضها
دون أن تتراكم رأسيًا. والنسبة $A/v$ هي ظلّ زاوية المسار
الواردة في القضية من الفصل~\ref{ch:b3:sequence-stratigraphy}: سمك كل
سحنة يساوي عرضها مضروبًا في ذلك الظلّ.
\end{solution}

\begin{solution}{exo:b3:facies-models:11}
الكثافة: $(2.65 - 2.30)/1.65 = 0.212$؛ \omterm{def:b3:facies-models:wireline}{السجلّ الصوتي}: $(86 - 52.6)/136.4 = 0.245$.
ومع $\phi = 0.245$، فإنّ $\rho_{\mathrm{ma}} = (2.30 - 0.245)/(1 - 0.245) =
\qty{2.72}{g/cm^3}$: أكثف من الكوارتز، كما في الحجر الرملي المتماسك بمادة
الكالسيت (2.71) أو الحاوي لحبيبات الكالسيت.
\end{solution}

\begin{solution}{exo:b3:facies-models:12}
الطين البحري المفتوح فوق الفحم يضع بيئات غير متجاورة في تماسّ: سطح
غمر، وهو هنا سطح تجاوزي ناتج عن تعرية الأمواج، مع فتات الأصداف
رواسب متخلّفة. ولمّا تراجع خط الشاطئ نحو اليابسة، قطعت الأمواج القاعَ
إلى نحو \omterm{def:b3:facies-models:shoreface}{قاعدة أمواج الطقس الهادئ}: الشاطئ، \omterm{def:b3:facies-models:shoreface}{والوجه الشاطئي} العلوي،
وقمة السهل الساحلي، حتى نحو \qty{10}{m} من المقطع، قد تكون
مفقودة.
\end{solution}

\begin{solution}{pb:b3:facies-models:1}
\textbf{1.} الطباشير (بارد، كثافة منخفضة)، والطَّفل من \qty{766}{m} (حارّ)، والحجر الجيري من
\qty{866}{m} (بارد، كثيف، سريع)، والحجر الرملي من \qty{1046}{m} (أشكال
الجرس)، والشيست من \qty{1166}{m}.

\textbf{2.} الطباشير \qty{766}{m}، والطَّفل 100، والحجر الجيري 180، والحجر الرملي
\qty{120}{m}؛ قاعدة الغطاء عند $441.0 - 1166 = \qty{-725}{m}$.

\textbf{3.} $(2.71 - 2.03)/1.71 = 0.40$ قرب السطح، $(2.71 -
2.36)/1.71 = 0.20$ عند \qty{750}{m}: يفقد الطباشيرُ نصف مسامّيته
تحت وزن الرواسب الفوقية، بالانضغاط.

\textbf{4.} حبيباته طينيات ذات كثافة غير مؤكَّدة، وماؤه مرتبط
جزئيًا بها: تفشل افتراضات المعادلة.

\textbf{5.} الأساس البلوري: يُخترق الغطاء، ولا يُنتظر وجود
أيّ مكمن تحته.

\textbf{6.} $(2.71 - 2.573)/1.71 = 0.080$.

\textbf{7.} $(58.9 - 47.6)/141.4 = 0.080$: النتيجة نفسها.

\textbf{8.} $(2.65 - 2.32)/1.65 = 0.200$.

\textbf{9.} $(79.9 - 52.6)/136.4 = 0.200$: النتيجة نفسها.

\textbf{10.} $0.33/1.55 = 0.213$، أي أكبر بنقطة وثلث.

\textbf{11.} الحجر الرملي، بمسامية \qty{20}{\percent} في رماله؛
الطَّفل فوق الحجر الجيري يسدّه، وطينيات السهل الفيضي تقسّم
الحجر الرملي إلى مقصورات.

\textbf{12.} شكل جرسي: أشعة غاما منخفضة في الرمل المتصالب الطبقات، ترتفع عبر
فترة التناعم حتى القيم العالية للطين المتجذّر، على قاعدة حادّة.

\textbf{13.} \omterm{def:b3:facies-models:association}{التعاقب المتناعم نحو الأعلى}: رواسب باقية في قاع القناة فوق قاعدة تعرية،
كثبان القناة (طبقات متصالبة)، قمّة \omterm{def:b3:facies-models:rivers}{الحاجز النقطي} (تموّجات صغيرة)، السهل الفيضي
(طين متجذّر).

\textbf{14.} القناة، وحاجزها النقطي على الضفة الداخلية، والجسور الطبيعية
والسهل الفيضي ذو التربة، جنبًا إلى جنب.

\textbf{15.} نحو سماكة الرمل وفترة التناعم، $10.4 + 3.2 =
\qty{13.6}{m}$، أكثر قليلًا قبل الانضغاط.

\textbf{16.} نهر متعرّج تحرّك حزام قناته مرارًا بفعل
\omterm{def:b3:facies-models:rivers}{انقلاب المجرى} بينما انخفض الحوض، فتقطع كل قناة جديدة في
السهل الفيضي الذي خلّفته القناة السابقة.

\textbf{17.} لا: كل قاعدة ناتجة عن التعرية؛ طين السهل الفيضي الذي تحتها ورمل
القناة الذي فوقها لا يلزم أن يكونا متجاورين.

\textbf{18.} $(60 - 30)/(105 - 30) = 0.40$.

\textbf{19.} الأدنى: $14\times(0.60 + 0.40\times0.25) = \qty{9.80}{m}$؛
الأكبر: $16\times(0.65 + 0.40\times0.20) = \qty{11.68}{m}$.

\textbf{20.} الرمل النقي \qty{58.8}{m}؛ فترات التناعم \qty{33.4}{m}؛ الطين
\qty{27.8}{m}.

\textbf{21.} الصافي $= 58.8 + 0.40\times33.4 = \qty{72.2}{m}$؛ \omterm{def:b3:facies-models:ntg}{نسبة الصافي إلى الإجمالي}
$72.2/120 = 0.60$.

\textbf{22.} الكسور 0.20 و0.60: \qty{65.5}{m} و\qty{78.8}{m}، بفارق
مقداره \qty{13}{m} ناتج عن اختيار حدّ القطع وحده.

\textbf{23.} $72.2\times0.20\times10^6 = \qty{1.44e7}{m^3}$ لكل
كيلومتر مربّع.

\textbf{24.} سماكة الرمل الصافية للحجر الرملي الترياسي عند فاردن-2 هي
\textbf{\qty[text-series-to-math=true]{72.2}{m}}.
\end{solution}
